\section*{Введение}

Объектное хранилище, совместимое с протоколом \gls{s3}, является одним из ключевых компонентов современной инфраструктуры: в нём размещаются пользовательский контент, медиафайлы, бэкапы и артефакты сборок.
В <<Авито>> доступ к такому хранилищу предоставляется внутренним командам через платформу \gls{staas} (\textit{Storage as a Service}), объединяющую десятки \gls{s3}-кластеров и петабайтные объёмы данных.

По мере роста платформы возникает потребность в систематическом переносе бакетов между кластерами.
Основными причинами миграций являются плановое обновление железа и вывод устаревших кластеров из эксплуатации, а также разбиение чрезмерно больших кластеров, негативно влияющих на производительность и время восстановления после отказов.
Это делает миграцию бакетов регулярной операцией платформы, а не разовым мероприятием.

В настоящее время она проводится вручную силами SRE-инженеров, и у такого подхода есть несколько серьёзных проблем:
\begin{itemize}
    \item высокие трудозатраты на подготовку, копирование данных и согласование переключения с командой-владельцем;
    \item вынужденный простой на этапе переключения, когда клиент приостанавливает запись, что напрямую сказывается на пользовательских сценариях;
    \item риск рассинхронизации источника и приёмника, ошибок в конфигурации доступов и неполного переноса объектов~--- например, из-за ошибки в скрипте миграции.
\end{itemize}
Цена этих проблем растёт вместе с размером бакета:
\begin{itemize}
    \item каждое неверное действие администратора обходится дороже;
    \item влияние на пользовательские сценарии становится заметнее.
\end{itemize}

Это и мотивирует разработку автоматизированного механизма миграции бакетов между \gls{s3}-кластерами и его интеграцию в платформу \gls{staas}.
На первом шаге создаётся MVP-сервис на базе \gls{rclone}: он снимает с администраторов рутину переноса, но всё ещё требует короткой остановки записи на момент переключения.
Готовый инструмент не позволяет избавиться от простоя полностью: он не умеет брать блокировки на объекты, поэтому при параллельной записи во время миграции возможны гонки и потеря данных.
По этой причине на втором шаге разрабатывается связка из собственного мигратора и прокси перед бакетом: мигратор копирует данные с поддержкой распределённых блокировок, а прокси прозрачно для клиента маршрутизирует запросы между исходным и целевым бакетом на время переноса.
